\section{No Fixed-Points Beyond Second Order}
\label{sec:nogo}

In this section, we sketch an argument showing that, in the present setting,
W-types beyond \AgdaFunction{2W} cannot exist. Suppose the general higher-order
W-types exist, that is, we \AgdaKeyword{postulate}

\begin{code}
\>[0]\AgdaPostulate{HW}\AgdaSpace{}%
\AgdaSymbol{:}\AgdaSpace{}%
\AgdaFunction{HCont}\AgdaSpace{}%
\AgdaSymbol{(}\AgdaGeneralizable{A}\AgdaSpace{}%
\AgdaOperator{\AgdaInductiveConstructor{⇒}}\AgdaSpace{}%
\AgdaGeneralizable{A}\AgdaSymbol{)}\AgdaSpace{}%
\AgdaSymbol{→}\AgdaSpace{}%
\AgdaFunction{HCont}\AgdaSpace{}%
\AgdaGeneralizable{A}\<%
\end{code}


Within the calculus, we can define the signature of the ordinary W-type as a
higher-order container:

\begin{code}
\>[0]\AgdaFunction{SigWᶜ}\AgdaSpace{}%
\AgdaSymbol{:}\AgdaSpace{}%
\AgdaFunction{HCont}\AgdaSpace{}%
\AgdaSymbol{(((}\AgdaInductiveConstructor{*}\AgdaSpace{}%
\AgdaOperator{\AgdaInductiveConstructor{⇒}}\AgdaSpace{}%
\AgdaInductiveConstructor{*}\AgdaSymbol{)}\AgdaSpace{}%
\AgdaOperator{\AgdaInductiveConstructor{⇒}}\AgdaSpace{}%
\AgdaInductiveConstructor{*}\AgdaSymbol{)}\AgdaSpace{}%
\AgdaOperator{\AgdaInductiveConstructor{⇒}}\AgdaSpace{}%
\AgdaSymbol{(}\AgdaInductiveConstructor{*}\AgdaSpace{}%
\AgdaOperator{\AgdaInductiveConstructor{⇒}}\AgdaSpace{}%
\AgdaInductiveConstructor{*}\AgdaSymbol{)}\AgdaSpace{}%
\AgdaOperator{\AgdaInductiveConstructor{⇒}}\AgdaSpace{}%
\AgdaInductiveConstructor{*}\AgdaSymbol{)}\<%
\\
\>[0]\AgdaFunction{SigWᶜ}\AgdaSpace{}%
\AgdaSymbol{=}\AgdaSpace{}%
\AgdaInductiveConstructor{lam}\AgdaSpace{}%
\AgdaSymbol{(}\AgdaInductiveConstructor{lam}\AgdaSpace{}%
\AgdaSymbol{(}\AgdaFunction{napp}\AgdaSpace{}%
\AgdaSymbol{(}\AgdaFunction{nvar}\AgdaSpace{}%
\AgdaInductiveConstructor{vz}\AgdaSymbol{)}\AgdaSpace{}%
\AgdaSymbol{(}\AgdaFunction{napp}\AgdaSpace{}%
\AgdaSymbol{(}\AgdaFunction{nvar}\AgdaSpace{}%
\AgdaSymbol{(}\AgdaInductiveConstructor{vs}\AgdaSpace{}%
\AgdaInductiveConstructor{vz}\AgdaSymbol{))}\AgdaSpace{}%
\AgdaSymbol{(}\AgdaFunction{nvar}\AgdaSpace{}%
\AgdaInductiveConstructor{vz}\AgdaSymbol{))))}\<%
\end{code}


morally, \AgdaFunction{SigWᶜ} represents the operation
\AgdaBound{F}\;\AgdaBound{G}\;\AgdaSymbol{↦}\;\AgdaBound{G}\;\AgdaSymbol{(}\AgdaBound{F}\;\AgdaBound{G}\AgdaSymbol{)}. Combining the two, we obtain a so-called \emph{internal} W-type. By internal, we mean an operator represented as a higher-order container in our calculus.

\begin{code}%
\>[0]\AgdaFunction{Wᶜ}\AgdaSpace{}%
\AgdaSymbol{:}\AgdaSpace{}%
\AgdaFunction{HCont}\AgdaSpace{}%
\AgdaSymbol{((}\AgdaInductiveConstructor{*}\AgdaSpace{}%
\AgdaOperator{\AgdaInductiveConstructor{⇒}}\AgdaSpace{}%
\AgdaInductiveConstructor{*}\AgdaSymbol{)}\AgdaSpace{}%
\AgdaOperator{\AgdaInductiveConstructor{⇒}}\AgdaSpace{}%
\AgdaInductiveConstructor{*}\AgdaSymbol{)}\<%
\\
\>[0]\AgdaFunction{Wᶜ}\AgdaSpace{}%
\AgdaSymbol{=}\AgdaSpace{}%
\AgdaPostulate{HW}\AgdaSpace{}%
\AgdaFunction{SigWᶜ}\<%
\end{code}


Correspondingly, the \emph{external} W-type means the functor derived by applying the
functorial interpretation on the internal one:

\begin{code}%
\>[0]\AgdaFunction{W}\AgdaSpace{}%
\AgdaSymbol{:}\AgdaSpace{}%
\AgdaFunction{HCont}\AgdaSpace{}%
\AgdaSymbol{(}\AgdaInductiveConstructor{*}\AgdaSpace{}%
\AgdaOperator{\AgdaInductiveConstructor{⇒}}\AgdaSpace{}%
\AgdaInductiveConstructor{*}\AgdaSymbol{)}\AgdaSpace{}%
\AgdaSymbol{→}\AgdaSpace{}%
\AgdaFunction{HCont}\AgdaSpace{}%
\AgdaInductiveConstructor{*}\<%
\\
\>[0]\AgdaFunction{W}\AgdaSpace{}%
\AgdaSymbol{=}\AgdaSpace{}%
\AgdaFunction{napp}\AgdaSpace{}%
\AgdaFunction{Wᶜ}\<%
\end{code}


Up to the equivalences
\AgdaFunction{HCont}\;\AgdaSymbol{(}\AgdaInductiveConstructor{*}\;\AgdaOperator{\AgdaInductiveConstructor{⇒}}\;\AgdaInductiveConstructor{*}\AgdaSymbol{)}\;\AgdaSymbol{≃}\;\AgdaRecord{Cont}
and
\AgdaFunction{HCont}\;\AgdaInductiveConstructor{*}\;\AgdaSymbol{≃}\;\AgdaPrimitive{Type},
the external operator \AgdaFunction{W} has type
\AgdaRecord{Cont}\;\AgdaSymbol{→}\;\AgdaPrimitive{Type} and coincides with the
ordinary W-type defined inductively using Agda \emph{data} keyword: \AgdaKeyword{data}\;\AgdaDatatype{W}\;\AgdaSymbol{(}\AgdaBound{SP}\;\AgdaSymbol{:}\;\AgdaPostulate{Cont}\AgdaSymbol{)}\;\AgdaSymbol{:}\;\AgdaPrimitive{Type}. This is a satisfying result so far, which means we can present external W-types inside our calculus.

However, if we instead use the type interpretation, we obtain

\begin{code}%
\>[0]\AgdaFunction{W-bad}\AgdaSpace{}%
\AgdaSymbol{:}\AgdaSpace{}%
\AgdaSymbol{(}\AgdaPrimitive{Type}\AgdaSpace{}%
\AgdaSymbol{→}\AgdaSpace{}%
\AgdaPrimitive{Type}\AgdaSymbol{)}\AgdaSpace{}%
\AgdaSymbol{→}\AgdaSpace{}%
\AgdaPrimitive{Type}\<%
\\
\>[0]\AgdaFunction{W-bad}\AgdaSpace{}%
\AgdaSymbol{=}\AgdaSpace{}%
\AgdaOperator{\AgdaFunction{H⟦}}\AgdaSpace{}%
\AgdaFunction{Wᶜ}\AgdaSpace{}%
\AgdaOperator{\AgdaFunction{⟧}}\<%
\end{code}  


which assigns a fixed-point
\AgdaFunction{W-bad}\;\AgdaBound{F}\;\AgdaSymbol{≃}\;\AgdaBound{F}\;\AgdaSymbol{(}\AgdaFunction{W-bad}\;\AgdaBound{F}\AgdaSymbol{)}
to \emph{every} function
\AgdaBound{F}\;\AgdaSymbol{:}\;\AgdaPrimitive{Type}\;\AgdaSymbol{→}\;\AgdaPrimitive{Type}.
This cannot hold, since fixed-points are only guaranteed to exist for strictly
positive functors: taking
\AgdaBound{F}\;\AgdaBound{X}\;\AgdaSymbol{=}\;\AgdaBound{X}\;\AgdaSymbol{→}\;\AgdaDatatype{⊥},
we would obtain a type \AgdaBound{D} with
\AgdaBound{D}\;\AgdaSymbol{≃}\;\AgdaSymbol{(}\AgdaBound{D}\;\AgdaSymbol{→}\;\AgdaDatatype{⊥}\AgdaSymbol{)},
from which \AgdaDatatype{⊥} is derivable. Hence we arrive at a contradiction.
